\paragraph{A finite simultaneous set of detecting elements.}
Let $G$ be profinite, $A=k[[t]]$ with $k$ finite of characteristic
$\ell$, and suppose every open subgroup $U$ has finitely many
continuous homomorphisms $U\to\mathbf F_\ell$.

\begin{lemma}[Finite generation and finite detection]
\label{lem:selmer-finite-selection}
For a continuous action on a finite free $A$-module $M$, the module
$Z=Z^1_{\mathrm{cont}}(G,M)$ is finite free. Any family of linear forms
$\lambda_s:Z\to A$ has a subfamily of at most $\operatorname{rank}_A Z$
members with exactly the same common kernel over $A$.
\end{lemma}

\begin{proof}
Let $U$ be the open kernel of the residual action. A cocycle into
$M/tM$ is determined by its homomorphism on $U$ and its values on
coset representatives. Thus the residual cocycle space is finite.
Continuous functions $G\to M$ are complete for uniform $t$-adic
convergence, and the cocycle equations are closed. Dividing a cocycle
with values in $tM$ by $t$ preserves continuity and the equations.
It follows that $Z$ is complete and
$Z/tZ\hookrightarrow Z^1_{\mathrm{cont}}(G,M/tM)$.
Lifting finitely many generators of $Z/tZ$ and successively correcting
the remainder modulo powers of $t$ generates $Z$ over $A$.
It is torsion-free, hence free.

For the second assertion, extend the forms to the fraction field and
choose a basis of their span from the forms themselves. Vanishing of
these basis forms forces every original form to vanish in the fraction
field, hence in $A$. No extra denominator enters the common kernel.
\end{proof}

The lemma also constructs the odd family used in integral detection.
The actual odd Selmer conditions define a submodule $S$ of the finite
cocycle module, because their local witnesses can be added and scaled.
Let $B$ be the skew odd boundaries in all cochains. Choose a
fraction-field basis of $(S/(S\cap B))\otimes_A k((t))$ from images of
actual elements $\kappa_i\in S$. They are integrally independent
modulo boundaries. Expressing any class in that basis, clearing
denominators and killing the remaining torsion gives generation up to
a power of $t$. The original local witnesses of each $\kappa_i$ are
retained. This also covers an empty family and does not require
$B\subset S$.

Apply the lemma at each of finitely many normalized points to the
original full matrix module over $D_A$, with action
$Y\mapsto\chi(g)\rho_1(g)Y\rho_1(g)^{-1}$, and the forms
\[
                  \lambda_h(z)=[\eta]\operatorname{tr}(z(h)\rho_1(h)).
\]
The actual paired cocycles form a submodule. Retain the spectral
hypothesis that the elements with separable characteristic polynomial
are dense in the original cyclotomic kernel $H$ at each point. Each
such locus is open: a nonzero discriminant is detected modulo some
power of the uniformizer. Their finite intersection $S\subset H$ is
dense. For fixed $z$, the displayed expression is continuous in $h$;
vanishing on $S$ therefore implies vanishing on $H$.

Proposition~\ref{prop:selmer-integral-paired-detection} and
Lemma~\ref{lem:selmer-finite-selection} now give a finite subset of $S$
detecting at each point. Their union is one finite set detecting at
every point, and all its characteristic polynomials are separable at
every point. Its size is at most the sum of the ranks of these
continuous cocycle modules \cite{alpo-finite-selection}.

There is also one denominator exponent for all integral cocycles in
the detecting kernel at a \emph{fixed} point. Indeed, that kernel is a
submodule of a finite module over $A$, hence has finitely many
generators. Choose the polarized primitives supplied by integral
detection for these generators, multiply them by powers of $t$ to
make their exponents equal to the maximum, and take linear
combinations. This preserves the original coefficient module and
polarization. The exponent depends on the fixed arithmetic data.
This argument neither lifts modular cocycles nor replaces the
separate count of all solutions modulo $t^a$. In particular, equality
of integral kernels need not persist modulo $t^a$: the forms $z$ and
$tz$ already illustrate this distinction.
